
\section{Diskussion}
Die prototypische Umsetzung zeigt, dass ein gestenbasiertes Marking-Menu als zusätzliche Interaktionsschicht für den Multisensorischen Interaktionsraum grund\-sätzlich realisierbar ist. Insbesondere die modulare Trennung zwischen Benutzungsoberfläche und Kommunikation mit der Zielanwendung erwies sich als vorteilhaft, da einzelne Komponenten unabhängig voneinander entwickelt und getestet werden konnten. Die angestrebte Szenariounabhängigkeit des Systems wird dabei weniger durch die reine Architektur als vielmehr durch die datengetriebene Menüstruktur sowie die abstrahierte Eventkommunikation unterstützt.

Eine zentrale Stärke des entwickelten Prototyps liegt daher in der flexiblen, datengetriebenen Menüstruktur. Da die Menükonfiguration nicht fest im Overlay verankert ist, sondern von der Zielanwendung bereitgestellt wird, lässt sich die Interaktionsoberfläche an unterschiedliche Anwendungskontexte anpassen, ohne Änderungen an der Interaktionslogik oder der Benutzungsoberfläche selbst vorzunehmen. In Kombination mit der WebSocket-basierten Kommunikation entsteht so eine lose Kopplung zwischen Interaktionssystem und Zielanwendung. Auch wenn die konkrete Umsetzung der Nachrichtenstruktur noch nicht vollständig konsistent ist, erwies sich dieser Ansatz grund\-sätzlich als geeignet, um Interaktionsereignisse in Echtzeit zu übertragen und unterschiedliche Anwendungsszenarien anzubinden.

Gleichzeitig zeigte die praktische Umsetzung mehrere technische und konzeptionelle Grenzen des Systems. Eine zentrale Herausforderung betrifft die Robustheit der Gestenerkennung. Die verwendeten Gestenheuristiken basieren auf einfachen Distanzmaßen zwischen Landmarken und sind dadurch zwar effizient und nachvollziehbar, jedoch anfällig gegenüber Perspektivwechseln, unvollständiger Erfassung sowie unterschiedlichen Abständen zur Kamera. Während diese Problematik in der Entwicklung auf einem Laptop kaum sichtbar war, trat sie im Multisensorischen Interaktionsraum deutlich hervor, da sich Nutzer:innen frei im Raum bewegen und somit in variierenden Distanzen und Winkeln zur Kamera befinden.

Insbesondere bei der Pinch-Geste zur Slidersteuerung kam es zu Fehlinterpretationen. Durch schräge Kameraperspektiven konnte es vorkommen, dass sich Daumen und Zeigefinger im Kamerabild scheinbar überlagerten, obwohl sie räumlich getrennt waren. Dies führte dazu, dass unbeabsichtigt eine Pinch-Geste erkannt wurde, beispielsweise beim Absenken der Hand, wodurch Parameterwerte unbeabsichtigt verändert wurden und Nutzer:innen explizit darauf achten mussten, Fehlinterpretationen zu vermeiden.

Auch die kontinuierliche Interaktion über Slider erwies sich als sensibler Bereich. Während die Steuerung grundlegender Parameter wie Lautstärke, Helligkeit oder Vibrationsintensität funktional zuverlässig funktionierte, zeigte sich insbesondere bei der Präsentationssteuerung ein Spannungsverhältnis zwischen Interaktivität und Systemreaktion. Bei längeren Videoinhalten oder häufigen Zustandsaktualisierungen kann die kontinuierliche Synchronisierung mit der Zielanwendung zu Verzögerungen führen. Gleichzeitig besteht bei Parametern wie Lautstärke oder Helligkeit die Erwartung, Änderungen unmittelbar und kontinuierlich wahrzunehmen. Daraus ergibt sich, dass unterschiedliche Interaktionstypen unterschiedliche Anforderungen an die Kommunikationsstrategie stellen: Während für kontinuierliche Parameter eine laufende Übertragung sinnvoll ist, kann für andere Anwendungsfälle eine Übermittlung erst nach Abschluss der Interaktion stabiler und effizienter sein.

Eine weitere Herausforderung zeigte sich im Mehrbildschirmbetrieb. Zwar konnte das System erfolgreich auf mehrere Projektionsflächen erweitert werden, jedoch erhöht diese räumliche Verteilung die Anforderungen an Kameraauswahl, Kalibrierung und Gestenerkennung erheblich. Die Einschränkung, Aktivierungsgesten nur bei ausreichend frontaler Handausrichtung zu erkennen, reduzierte Fehlaktivierungen zwischen mehreren Leinwänden, erwies sich jedoch nicht als vollständig robust. Die Zuordnung zwischen Nutzer:innen, Kameraansicht und Interaktionsfläche bleibt damit ein kritischer Bestandteil des Systems, insbesondere in dynamischen Nutzungssituationen.

Zusätzlich wurde deutlich, dass das System derzeit nicht für Mehrnutzerinteraktion ausgelegt ist. Es wird nicht unterschieden, von welcher Person oder welcher Hand eine Interaktion ausgeht. Dies führt zu Einschränkungen, etwa wenn das Menü ausgeblendet werden soll, da hierfür sichergestellt werden muss, dass gar keine Hand mehr im Kamerabild erkannt wird.

Auch die Stabilität der Cursorbewegung stellte eine Herausforderung dar. Im praktischen Einsatz zeigte sich ein deutliches Jitter in der Handerkennung, wodurch der Cursor unruhig über die Oberfläche springt. Dieses Verhalten trat insbesondere im MIR auf und war in der Laptop-basierten Entwicklungsumgebung kaum vorhanden. In der Folge kann es dazu kommen, dass der Cursor unbeabsichtigt Menügrenzen überschreitet, wodurch beispielsweise Submenüs geschlossen werden, obwohl die Interaktion fortgesetzt werden sollte. Eine mögliche Verbesserung wäre hier die Einführung erweiterter Interaktionsbereiche oder toleranter Kollisionszonen, die kurzfristige Positionssprünge abfangen und so die Interaktion stabilisieren.

Neben den genannten Herausforderungen zeigt sich auch auf architektonischer Ebene weiteres Verbesserungspotenzial. In der aktuellen Umsetzung sind die unterstützten Interaktionstypen bislang auf eine begrenzte Menge vordefinierter Aktionen beschränkt, insbesondere auf diskrete Auswahlaktionen sowie kontinuierliche Steuerung über Slider. Daraus ergibt sich eine teilweise noch vorhandene Kopplung zwischen der Darstellungsebene und der konkreten Bedeutung einzelner Aktionen. Langfristig erscheint es daher sinnvoll, die Interaktionslogik und die Nachrichtenstruktur für die Kommunikation weiter zu abstrahieren, sodass die Benutzungsoberfläche keine Kenntnis über die semantische Bedeutung einer Aktion besitzen muss. Stattdessen könnten Menüeinträge ausschließlich anhand ihres Interaktionstyps (zum Beispiel Button oder Slider) interpretiert werden, während die konkrete Funktionalität vollständig durch die Zielanwendung definiert wird. Die Benutzungsoberfläche würde in diesem Fall lediglich generische Interaktionsereignisse mit Referenzen auf das jeweilige Zielelement übermitteln. Ein solcher Ansatz würde die Entkopplung zwischen Interaktionssystem und Zielanwendung weiter stärken und die Erweiterbarkeit des Systems deutlich verbessern, da neue Aktionen ohne Anpassung der Benutzungsoberfläche integriert werden könnten.


Diese noch nicht vollständige Abstraktion zeigt sich insbesondere in der Nachrichtenstruktur.
Diese entstand im Rahmen eines iterativen Entwicklungsprozesses, bei dem zunächst einfache Interaktionsereignisse übertragen wurden und das Schema schrittweise erweitert wurde, um zusätzliche Anforderungen wie Initialisierung, Zustandssynchronisation und unterschiedliche UI-Elemente abzudecken. Dadurch entstand eine funktionale, jedoch nicht vollständig konsistente Struktur.
Beispielsweise besteht eine teilweise Redundanz zwischen den Feldern \texttt{``action''} und \texttt{``type''}, da bestimmte Kombinationen implizit festgelegt sind, etwa \texttt{``pressed''} nur für \texttt{``button''}. Ebenso ist das Feld \texttt{``value''} nicht strikt typisiert, da es sowohl kontinuierliche Werte als auch diskrete Zustände, beispielsweise \texttt{``play''}/\texttt{``pause''}, oder Zeitangaben enthalten kann.
Eine mögliche Weiterentwicklung des Protokolls könnte daher in einer stärkeren Trennung von Ereignistyp und Datenstruktur sowie einer strengeren Typisierung der übertragenen Werte liegen. Beispielsweise könnten unterschiedliche Nachrichtentypen oder spezialisierte Payload-Strukturen definiert werden, um die Semantik klarer abzubilden und die Erweiterbarkeit zu verbessern.

Darüber hinaus ist der aktuelle Prototyp bewusst als explorative Umsetzung zu verstehen und nicht als vollständig ausgereiftes Produktivsystem. Mehrere Aspekte wurden pragmatisch im Rahmen eines iterativen Entwicklungsprozesses realisiert. Dazu zählen die heuristische Gestenerkennung, die evolutionär gewachsene Nachrichtenstruktur sowie die Kalibrierung von Menü und Kameras. Diese Entscheidungen ermöglichten eine schnelle prototypische Umsetzung und eine frühzeitige Erprobung des Interaktionskonzepts, gehen jedoch zulasten von Standardisierung und technischer Robustheit.

Diese explorative Vorgehensweise spiegelt sich auch in der Art der durchgeführten Tests wider. Das Testen des Systems erfolgte iterativ und begleitend zur Entwicklung in unterschiedlichen Umgebungen. Während grundlegende Funktionalitäten fortlaufend in einer Laptop-basierten Entwicklungsumgebung getestet wurden, fanden insgesamt vier Tests im Multisensorischen Interaktionsraum zu unterschiedlichen Entwicklungszeitpunkten statt.
Diese Tests wurden in Zusammenarbeit mit einem Entwickler des MIR (Jimmy Orawetz) sowie dem Betreuer der Arbeit (Prof. Dr. Dietrich Kammer) durchgeführt. Die Testphasen umfassten sowohl frühe Entwicklungsstadien, in denen grundlegende Interaktionsmechanismen überprüft wurden, als auch fortgeschrittene Zustände des Systems, etwa nach der Integration des Overlays und der Zuordnung zu mehreren Projektionsflächen sowie abschließend nach der vollständigen Einbindung von Kameras und Interaktionslogik im Raum.
Die Rückmeldungen aus diesen Testphasen flossen direkt in die Weiterentwicklung des Systems ein und ermöglichten eine praxisnahe Bewertung technischer und konzeptioneller Entscheidungen. Auch wenn keine systematische empirische Evaluation durchgeführt wurde, kann diese Form der wiederholten, begleiteten Erprobung als informelle expertenbasierte Evaluation verstanden werden.

Trotz dieser Einschränkungen zeigt die praktische Umsetzung, dass das entwickelte Interaktionskonzept die im Abschnitt \ref{sec:synthese} formulierten Anforderungen grund\-sätzlich adressieren kann. Das System ermöglicht eine controllerfreie Interaktion, erhält die Bewegungsfreiheit im Raum und bietet über ein strukturiertes Menü einen klaren Zugriff auf systemweite Funktionen. Besonders hervorzuheben ist, dass menübasierte gestenbasierte Interaktion im MIR nicht nur konzeptionell, sondern auch technisch als eigenständige Interaktionsschicht realisierbar ist.

Für zukünftige Arbeiten ergeben sich daraus mehrere Entwicklungsperspektiven. Dazu zählen insbesondere robustere Verfahren zur Gestenerkennung, etwa durch lernbasierte Ansätze oder durch Algorithmen, die neben der Distanz einzelner Landmarken auch Relationen zur Handgröße sowie perspektivische Verzerrungen berücksichtigen. Ebenso erscheint eine differenziertere Behandlung kontinuierlicher Interaktionen sinnvoll, beispielsweise durch angepasste Synchronisationsstrategien oder spezialisierte Interaktionsmechanismen für unterschiedliche Parametertypen. Darüber hinaus könnte die Integration weiterer Modalitäten, etwa auditiver oder haptischer Rückmeldungen, das Interaktionssystem besser an den multisensorischen Charakter des MIR anpassen. Schließlich wäre eine systematische empirische Evaluation erforderlich, um Usability, Erlernbarkeit, körperliche Belastung und wahrgenommene Immersion im Nutzungskontext des MIR fundiert zu untersuchen.


\section{Zusammenfassung und Ausblick}
Ziel dieser Arbeit war die Konzeption und prototypische Umsetzung eines controllerfreien, körper\-basierten Interaktionssystems für den Multisensorischen Interaktionsraum (MIR). Im Fokus stand die Entwicklung eines auf Basis der Anwendungsanalyse identifizierten Interaktionskonzepts, das die zuvor abgeleiteten Anforderungen adressiert: eine Nutzung ohne zusätzliche Eingabegeräte, die Unterstützung freier Bewegung im Raum, ein geringer Integrationsaufwand sowie ein schneller und intuitiver Zugriff auf Interaktionsmöglichkeiten.

Auf Grundlage der Analyse bestehender Szenarien sowie der Auseinandersetzung mit verwandten Arbeiten wurde ein gestenbasiertes Marking-Menu als zentrales Interaktionspotenzial ausgewählt. Im Vergleich zu den ebenfalls betrachteten objektbasierten und positionsbasierten Interaktionsansätzen zeigt sich, dass der menübasierte Ansatz insbesondere durch seine Anwendungsunabhängigkeit und geringe Anforderungen an die Zielanwendung eine besonders hohe Übertragbarkeit auf unterschiedliche MIR-Szenarien besitzt. Dieses Konzept wurde daher im weiteren Verlauf konzeptionell ausgearbeitet und in Form eines prototypischen Interaktionssystems umgesetzt. Die entwickelte Lösung kombiniert eine radiale Menüstruktur mit gestenbasierter Auswahl, kontinuierlichen Steuerungsmechanismen und wird als eigenständige Interaktionsschicht über bestehende Anwendungen gelegt.

Die praktische Umsetzung zeigt, dass ein solches System grund\-sätzlich realisierbar ist und zentrale Anforderungen an körper\-basierte Interaktion im MIR adressieren kann. Insbesondere die modulare Architektur, die datengetriebene Menüstruktur sowie die lose Kopplung zwischen Interaktionssystem und Zielanwendung ermög\-lichen eine flexible Integration in unterschiedliche Anwendungsszenarien. Gleichzeitig verdeutlichen die Ergebnisse, dass die Gestaltung robuster und präziser Mid-Air-Interaktion weiterhin eine Herausforderung darstellt. Aspekte wie Gestenerkennung, kontinuierliche Steuerung und Mehrbildschirminteraktion erfordern eine sorgfältige Abstimmung zwischen technischer Umsetzung und Interaktionsdesign.

Die Arbeit leistet damit einen Beitrag zur Gestaltung controllerfreier Interaktion in projektionsbasierten, multisensorischen Umgebungen. Sie zeigt, wie menübasierte Interaktionskonzepte, insbesondere Marking-Menus, in einen räumlichen Kontext übertragen und mit gestenbasierter Eingabe kombiniert werden können. Darüber hinaus wird deutlich, dass die Integration kontinuierlicher Steuerungsmechanismen in solche Systeme eine bislang wenig untersuchte, aber relevante Erweiterung darstellt.

Für zukünftige Arbeiten ergeben sich mehrere Ansatzpunkte. Ein zentraler Aspekt ist die Verbesserung der Gestenerkennung, etwa durch robustere Algorithmen, die perspektivische Verzerrungen und unterschiedliche Abstände zur Kamera besser berücksichtigen. Ebenso erscheint eine differenzierte Behandlung kontinuierlicher Interaktionen sinnvoll, beispielsweise durch adaptive Synchronisationsstrategien.

Darüber hinaus stellt die Erweiterung des Systems um multimodale Interaktionsformen ein vielversprechendes Forschungsfeld dar. Die Kombination von Gesten mit auditivem oder haptischem Feedback könnte die Interaktion stärker in den multisensorischen Charakter des MIR einbetten und sowohl die Wahrnehmung als auch die Kontrolle verbessern.

Ein weiterer wichtiger Schritt besteht in der systematischen Evaluation des entwickelten Interaktionssystems. Empirische Studien könnten Aufschluss darüber geben, wie Nutzer:innen das System hinsichtlich Usability, Erlernbarkeit, körperlicher Belastung und Immersion wahrnehmen. Insbesondere der Vergleich unterschiedlicher Interaktionsmechanismen sowie die Untersuchung von Nutzungsszenarien mit mehreren Personen stellen hierbei relevante Fragestellungen dar.

Zusammenfassend zeigt die Arbeit, dass körper\-basierte, gestenbasierte Interaktion im Multisensorischen Interaktionsraum nicht nur konzeptionell, sondern auch technisch umsetzbar ist. Gleichzeitig wird deutlich, dass die Entwicklung robuster, intuitiver und kontextangemessener Interaktionssysteme in solchen Umgebungen weiterhin ein offenes und vielschichtiges Forschungsfeld darstellt.
